% =====================================================================
%  CHAPTER 1 — INTRODUCTION
% =====================================================================
\chapter{INTRODUCTION}
\label{ch:intro}

\section{Background and Motivation}
\label{sec:motivation}
Modern society depends heavily on space and radio systems: satellite navigation
(GNSS/GPS), radio communication, aviation, power grids, and Earth-observation
satellites. Most of these systems assume that the near-Earth space environment
stays reasonably stable. When the Sun becomes more active, \emph{space weather
driven by solar eruptive activity}-flares, coronal mass ejections (CMEs), and
the geomagnetic storms they cause-happens more often and can be more severe.
These events can disturb satellites, radio and GNSS links, and power grids. In
October 2024, NASA and the NOAA Space Weather Prediction Center jointly
announced that Solar Cycle 25 had reached its maximum
phase~\citep{nasa_noaa_solarmax2024}. This means solar activity, and its
practical effects, are especially relevant right now.

A clear example of this impact is the loss of 38 of the 49 Starlink satellites
launched by SpaceX in February 2022. A geomagnetic storm arrived while the
satellites were still at low altitude. It heated and expanded the thermosphere,
which increased atmospheric drag on the satellites much more than expected, and
most of them re-entered the atmosphere within days~\citep{berger2023starlink}.
This event shows that space weather is not only a scientific topic; it can
cause real economic and operational loss. Solar radio bursts (SRBs), which are
the radio-frequency signal produced by the same eruptive processes, can also
interfere directly with GNSS. A study of the International GNSS Service network
over Solar Cycle 24 found that SRBs can lower the carrier-to-noise density
ratio ($C/N_0$) of GPS and Galileo receivers, and in the worst cases reduce the
number of usable satellites at a station~\citep{floressoriano2024gnss}. This
links solar radio emission directly to a radio-frequency-interference problem,
and it is one reason to keep monitoring solar radio activity from the ground.

To understand and eventually forecast these effects, we need to observe the Sun
directly, and optical or thermal images alone are not enough: they show the
visible surface of the Sun but not the processes that actually drive space
weather. At radio wavelengths, the Sun is not only an optical or thermal
object; it is also a source that carries information about the coronal plasma,
energetic particles, and eruptive processes. Radio emission at metre and
decimetre wavelengths carries information about coronal plasma density and
temperature, shock fronts, and the energetic electron beams produced during
flares and CMEs~\citep{dulk1985,kundu1965,wild1963}. Because of this, observing
the Sun at radio frequencies is one of the oldest parts of radio astronomy, and
it remains important for space-weather work, since solar radio bursts are an
early, ground-observable sign of the eruptive events described above.

However, observing the Sun well at radio wavelengths usually means using large
and expensive equipment. Angular resolution depends on the ratio between the
observing wavelength and the antenna diameter. To resolve detail on the scale
of the \SI{32}{\arcminute} solar disk, a single dish would need to be tens of
metres across, together with the receivers, structure, and site infrastructure
that such a dish requires. This is difficult to build for most university
research groups, including those in Vietnam. \emph{Radio
interferometry}~\citep{tms2017,taylor1999synthesis} offers another way.
Instead of one large dish, \emph{aperture synthesis} combines the signals from
several small antennas placed apart from each other. The achievable resolution
then depends on the largest distance between antennas (the \emph{baseline}),
not on the size of a single antenna. This is the idea behind synthesis arrays
such as the historical Culgoora radioheliograph~\citep{mclean1985}, and modern
arrays like LOFAR~\citep{lofar2013} and the Murchison Widefield
Array~\citep{mwa2013}. It is also a reasonable direction for a low-cost solar
radio instrument.

This low-cost, distributed approach is not just an idea used in this thesis; it
is already an active research direction. DLITE (the Deployable Low-band
Ionosphere and Transient Experiment) is a four-element, field-deployable
interferometer that works at \SIrange{30}{40}{\mega\hertz}, built from four
commercially available LWA-style crossed-dipole antennas~\citep{carson2022dlite}.
The published total system cost of DLITE, including receivers, digitisers, and
the rest of the back end, is around US\$45{,}000~\citep{carson2022dlite}, but
this figure is not what makes the antenna itself low-cost: official LWA
component pricing lists the crossed-dipole antenna at about \$793, its
front-end-electronics balun at about \$308, and its ground stake at about
\$65.50~\citep{lwa_pricing}, so the raw receiving element used by DLITE costs
on the order of \$1{,}200 per antenna, only a few thousand dollars for all four.
It is this element-level cost, not the fully-equipped system cost, that makes
the comparison with a single, tens-of-metres-across dish (Section~1.1)
meaningful. DLITE was built mainly to study the ionosphere, but it has also
detected long solar radio bursts, including Type~II and Type~IV events,
confirmed by independent e-Callisto stations~\citep{carson2022dlite}. This
shows that a small, distributed, low-cost array is a realistic way to monitor
solar radio bursts, not only something for large professional facilities.

This thesis works at \SI{610}{\mega\hertz}, in the UHF/decimetric radio band.
This is a practical choice, not the frequency where the strongest solar radio
bursts are usually found. Three reasons support it. First, Learmonth/SWS
already publishes daily solar flux data at \SI{610}{\mega\hertz}, so the
simulation can use a real, continuous, and freely available dataset instead of
only a synthetic source. Second, UHF Yagi-Uda antennas at this frequency are
cheap, easy to buy, and easy to build, which fits a low-cost, university-level
interferometer. Third, the wavelength at \SI{610}{\mega\hertz} is about
\SI{0.49}{\metre}, so a baseline of only a few metres to a few tens of metres
is already enough to produce a clear fringe pattern; this is convenient for
testing fringe period, visibility, $uv$-coverage, and dirty-image formation on
a small, deployable array. At the same time, \SI{610}{\mega\hertz} does not
represent the full picture of solar radio bursts: the strongest Type~II and
Type~III bursts are usually observed at lower frequencies (tens of MHz),
closer to the DLITE band described above. The choice of \SI{610}{\mega\hertz}
in this thesis reflects the available real data, the low antenna cost, and the
practical UHF direction of the project, not a claim that this band captures
every kind of solar radio burst.

Building an array like this, however, is harder than just placing several
antennas in a field. A working interferometer needs receiver channels that stay
coherent in phase and clock with each other, a receiver noise level that is
well understood so it can be separated from the weak solar signal, an antenna
radiation pattern that is actually measured or modelled instead of assumed,
control of local radio-frequency interference (RFI), and an array layout that
gives good $uv$-coverage. These are not just theoretical concerns. In my own
preliminary two-element Yagi-Uda deployment at Xuan~Mai, Hanoi
(Section~\ref{subsec:field}), the recorded data showed only an ambiguous daily
trend, not a clean interferometric detection. This was traced to several
causes: phase-incoherent receiver oscillators, several decibels of cable
insertion loss, strong local RFI, and a receiver noise floor close to the
level of the quiet-Sun signal. Because of this experience, before spending
more effort on hardware for a larger low-cost solar radio interferometer, it
makes sense to first study and separate these effects in a controlled and
repeatable way.

This is what a physically grounded, end-to-end \emph{simulation} can provide,
and it is the subject of this thesis. Instead of treating the antenna, the
receiver noise, the array geometry, and the imaging step as separate problems,
the pipeline built in this thesis connects all of them using one real
observation: archived solar flux from the Learmonth Solar Radio Observatory
(Space Weather Services, Bureau of Meteorology,
Australia)~\citep{sws,learmonth_coords}. The flux is converted to a received
voltage using the Friis transmission formula and Johnson-Nyquist thermal
noise~\citep{friis1946,nyquist1928}, weighted by a real Yagi-Uda radiation
pattern modelled in 4NEC2~\citep{nec2} instead of a simplified Gaussian beam,
propagated through an array geometry to form visibilities, and then correlated
and inverted into a dirty image and point-spread function. All of this can be
checked before an array layout is actually built. Existing open-source
teaching tools such as \textsc{APSYNSIM}~\citep{apsynsim}, developed at the
Onsala Space Observatory for the Nordic ALMA Regional Centre Interferometry
School, give interactive visualisations of aperture-synthesis concepts, but
they use synthetic sky models and simplified Gaussian primary beams, without a
connection to real flux data, realistic noise, or an actual hardware path.
This thesis compares its simulation directly against \textsc{APSYNSIM}, to
show where a pipeline based on real data and real hardware adds value compared
to a purely educational tool, and also compares it against the field-deployed
two-element analogue interferometer at Xuan~Mai, to connect the simulation to
a real, if imperfect, measurement. Building and validating this end-to-end
simulation-so that array layouts, receiver budgets, and expected signal levels
can be checked before further field deployment-is the goal of this work.

\section{Problem Statement}
\label{sec:problem-statement}
The discussion above points to a gap, not just a list of separate technical
problems. A future low-cost, distributed UHF solar radio interferometer, of
the kind described in Section~1.1, should not be designed only by trial and
error in the field. As the Xuan~Mai deployment showed, an ambiguous result
from real hardware can come from several different causes at once:
phase-incoherent receivers, cable loss, local RFI, a weak array geometry, or a
genuinely weak signal. Without a reference model, it is hard to tell these
causes apart afterwards. The individual physics is already well understood:
how an antenna receives a flux density, how thermal receiver noise is added,
how the geometry of a moving baseline behaves, and how an image is formed from
correlated visibilities. What is missing is a single, physically grounded
software model that connects all of these stages using real data, so that the
behaviour of a proposed instrument can be predicted and checked \emph{before}
it is built. Without such a model, there is no reliable way to decide, before
a field deployment, whether a given receiver noise budget, antenna choice, or
array layout is good enough, or to tell an instrument fault apart from a real
design limitation. Closing this gap through simulation, rather than through
more hardware deployment, is the next necessary step, and it is the problem
this thesis addresses:

\begin{quote}\itshape
How can a physically grounded, end-to-end simulation be constructed to
evaluate the feasibility, signal chain, and array-geometry choices of a
future low-cost UHF solar radio interferometer before hardware deployment?
\end{quote}

\section{Objectives}
\label{sec:objectives}
To answer this question, the thesis pursues four objectives, each covering one
part of the simulated signal chain described above:
\begin{enumerate}[leftmargin=1.8em]
  \item \textbf{Antenna modelling.} Build a Method-of-Moments model of the
        commercial UHF Yagi-Uda receiving element in 4NEC2, and extract its
        input impedance, gain, and radiation pattern. Cross-check the result
        against a direct vector-network-analyser measurement and an
        independent finite-element solution. Also measure how much mutual
        coupling exists between elements at different baselines, and confirm
        that it is small enough to ignore at the separations used in this
        thesis.
  \item \textbf{End-to-end signal pipeline.} Build a Python pipeline that
        converts real SWS Learmonth \SI{610}{\mega\hertz} flux data into
        complex antenna voltages, using the Friis reception formula and
        Johnson-Nyquist noise. The pipeline then applies the per-antenna
        geometric phase, cross-correlates the signals to form visibilities,
        and inverts them into a dirty image.
  \item \textbf{General array geometry.} Replace any assumption of a
        one-dimensional, east-west-only array with a general
        local-to-equatorial ENU$\rightarrow$XYZ$\rightarrow\uvw$
        transformation, so that the simulation works for arbitrary
        two-dimensional array layouts, not just one fixed geometry.
  \item \textbf{Validation.} Check that the simulated instrument reproduces
        the expected analytic results-the fringe period, the $\jinc$-shaped
        visibility of a resolved disk, and the dirty image as the sky
        convolved with the synthesised beam-using the transit of
        \mbox{20 March 2025} at Learmonth.
\end{enumerate}

\section{Contributions}
\label{sec:contributions-intro}
The main contributions of this thesis are:
\begin{enumerate}[leftmargin=1.8em]
  \item A 4NEC2 Method-of-Moments model of a low-cost commercial UHF Yagi-Uda
        element, cross-checked against direct measurement and an independent
        finite-element solution, together with a mutual-coupling study
        showing that the elements are effectively independent at the
        operating separations, with coupling below $-78~\si{dB}$ beyond
        \SI{5}{\metre}.
  \item A modular, reproducible interferometry simulation engine that works
        for \emph{arbitrary} two-dimensional array geometries, using one
        consistent $\uvw$ formulation for per-antenna phase, baseline
        computation, and fringe stopping.
  \item A demonstration that the simulated instrument, driven by real
        archived solar flux, reproduces the classical analytic results of
        aperture synthesis-directly answering the feasibility question
        raised in Section~1.2, and validating both the instrument concept
        and the simulation itself as a design tool.
  \item A structured comparison between this simulation and
        \textsc{APSYNSIM}~\citep{apsynsim}, the most widely used open-source
        aperture-synthesis educational simulator, showing where a pipeline
        based on real data and real hardware adds value compared to a purely
        educational tool.
  \item A reproducible, single-notebook workflow that takes a new user from
        an archived SWS daily file to a calibrated dirty image in under two
        minutes of computation, with all parameters set from a single
        configuration cell.
\end{enumerate}

\section{Organization of the Thesis}
The rest of this thesis has five chapters.
\textbf{Chapter~2} covers the theoretical background: solar radio emission,
the basics of antenna theory, receiver noise, the principle of interferometry
and the van Cittert-Zernike theorem, and aperture-synthesis imaging.
\textbf{Chapter~3} presents the electromagnetic modelling of the Yagi-Uda
receiving element and the mutual-coupling study of the array.
\textbf{Chapter~4} describes the design of the end-to-end simulation
pipeline, from flux-data acquisition through voltage synthesis, array
geometry, correlation, and imaging.
\textbf{Chapter~5} reports and analyses the simulation results, including
validation against analytic predictions and the 20 March 2025 transit.
The final chapter, \textbf{Conclusion and Future Work}, summarises the
findings and outlines possible extensions. References and appendices follow.
\clearpage
